%!TEX root = ../main.tex

\section{Einleitung}

% Prüfzifferverfahren sind ein essenzielles Werkzeug zur Sicherstellung der Datenintegrität in zahlreichen Anwendungen des täglichen Lebens. Sie dienen dazu, Eingabefehler, wie etwa Tippfehler oder Übertragungsfehler, zu erkennen und somit die Zuverlässigkeit von Daten zu erhöhen.

% In der modernen Informationsverarbeitung spielen Prüfziffern eine zentrale Rolle, beispielsweise bei der Validierung von Bankkontonummern, ISBN-Nummern für Bücher oder internationalen Standardcodes wie der IBAN. Diese Verfahren basieren auf mathematischen Prinzipien der Diskreten Mathematik und ermöglichen eine effiziente Fehlererkennung.

% Im Folgenden werden die Grundlagen der Prüfzifferverfahren erläutert, verschiedene Methoden vorgestellt und deren praktische Anwendungen diskutiert. Ziel ist es, ein tieferes Verständnis für die Funktionsweise und die Bedeutung dieser Verfahren zu vermitteln.
\chapter{Prüfzifferverfahren}


\section{ISBN}
%!TEX root = ../main.tex

\section{EAN}
%!TEX root = ../main.tex

\section{EKONS}

%!TEX root = ../main.tex

\chapter{Fazit}
Aufgabenstellung, Vorgehensweise und wesentliche Ergebnisse werden kurz und präzise dargestellt und kritisch reflektiert. Die Zusammenfassung ist eigenständig verständlich. Länge ca. 1 bis 1,5 Seiten (Problem, Ziele, Vorgehensweise, Ergebnisse und Ausblick).

\chapter{Ausblick}


%% Beispiele, können problemlos entfernt werden!
\chapter{Beispiele}
Beispiele, können problemlos entfernt werden!
\section{Literatur}
Beispiel Text \cite[S. 10]{testBuch1}
Beispiel Homepage\cite{urlId} \\

\section{Bilder}
\begin{figure}[H]
	\centering
	\includegraphics[width=0.3\linewidth]{resources/images/logo-dhbw}
	\caption{DHBW Logo}
	\label{fig:logo-dhbw}
\end{figure}

\section{Fußnote und Abkürzung}
Fußnote\footnote{Fußnote}, \ac{DHBW}

\section{Tabelle}
\begin{table}[H]
	\centering
	\caption{Tabellenbeispiel}
	\label{tab:example}
	\begin{tabular}{|l|c|r|}
		\hline
		Spalte 1 & Spalte 2 & Spalte 3 \\
		\hline
		Zeile  &  &  \\
		\hline
		& Zeile &  \\
		\hline
		&  & Zeile \\
		\hline
		\multicolumn{2}{|r|}{Verbunden}	& nicht Verbunden \\
		\hline
	\end{tabular}
\end{table}

\section{Skript}
\lstinline[language=c]|printf("Inline Code");|
\lstinputlisting[language=python,caption={Beispiel python Skript},captionpos=t,label=scr:exanple]{resources/example-script.py}

\section{TODOs}
Text Text\todo{TODO als Randnotiz} Text
\todo[inline]{TODO als in Zeile}

\begin{landscape}
	\section{Querformat}
	Text über Tabelle \ref{tab:breitetabelle}.
	\begin{table}[H]
		\caption{Breites Tabellenbeispiel}
		\label{tab:breitetabelle}
		\centering
		\resizebox{\columnwidth}{!}{%
			\begin{tabular}{|p{10cm}|p{10cm}|p{10cm}|}
				\hline
				Sehr & breite & Tabelle \\
				\hline
				\multicolumn{3}{|c|}{Tabelle wurde in der Größe verkleinert} \\
				\hline
			\end{tabular}
		}
	\end{table}
	Text unter Tabelle \ref{tab:breitetabelle}.
\end{landscape}

\section{Auflistungen}
\begin{enumerate}
	\item Aufzählung nummeriert
\end{enumerate}
\begin{itemize}
	\item Aufzählung Stichpunkte
\end{itemize}
\begin{description}[style=nextline]
	\item[Label] Aufzählung als Beschreibung
\end{description}
